% Part: incompleteness
% Chapter: introduction
% Section: decidability

\documentclass[../../../include/open-logic-section]{subfiles}

\begin{document}

\olfileid{inc}{int}{dec}

\olsection{Ora Bisa Diputusake lan Ora Jangkep}

Buktine G\"odel kanggo teorema-teorema ora
jangkep mbutuhake aritmetisasi sintaks. Nanging
tanpa iku uga kita bisa entuk sawetara asil
kang migunani, mung kanthi nganggep yen sawijining
teori !!{represents} kabeh relasi !!{decidable}.
Buktine nggunakake argumen diagonal kang padha
karo bukti yen masalah mandheg ora bisa diputusake.

\begin{thm}
Yen $\Gamma$ teori konsisten kang !!{represents}
saben relasi !!{decidable}, mula $\Gamma$
ora !!{decidable}.
\end{thm}

\begin{proof}
Upamane $\Gamma$ iku !!{decidable}.
Kita bakal nuduhake yen $\Gamma$
!!{represents} saben relasi !!{decidable},
teori iku mesthi ora konsisten.

Sifat-sifat !!^{decidable} (relasi siji papan)
digambarake dening !!{formula}s kanthi siji
variabel bebas. Tetepna $!A_0(x)$, $!A_1(x)$,
\dots, minangka enumerasi komputabel kabeh
!!{formula}s kaya mangkono. Saiki gatekna
himpunan $D \subseteq \Nat$ ing ngisor iki:
\[
D = \Setabs{n}{\Gamma \Proves \lnot !A_n(\num{n})}
\]
Himpunan $D$ iku !!{decidable}, amarga kita
bisa mriksa apa $n \in D$ kanthi ngitung
$!A_n(x)$ dhisik, banjur saka iku
$\lnot !A_n(\num{n})$. Mesthi wae,
ngganti nganggo suku $\num{n}$ saben
kemunculan bebas $x$ ing $!A_n(x)$,
banjur menehi $!A(\num{n})$ awalan $\lnot$
bisa ditindakake kanthi mekanis. Miturut
asumsi, $\Gamma$ iku !!{decidable}, mula
kita bisa mriksa apa
% OLPL-168: Sumber ing langkah tes iki nulis A tanpa indeks n, senajan definisi D lan konstruksi sadurunge nggunakake A_n.
$\lnot !A(\num{n}) \in \Gamma$.
Yen iya, $n \in D$; yen ora,
$n \notin D$. Mula $D$ uga !!{decidable}.

Amarga $\Gamma$ !!{represents} kabeh sifat
!!{decidable}, teori iku !!{represents}~$D$.
Lan kabeh !!{formula}s kang !!{represents}s $D$
ing~$\Gamma$ ana ing antarane $!A_0(x)$,
$!A_1(x)$, \dots. Mula tetepna $d$ minangka
bilangan saengga $!A_d(x)$ !!{represents} $D$
ing~$\Gamma$. Yen $d \notin D$, mula amarga
$!A_d(x)$ !!{represents}~$D$,
$\Gamma \Proves \lnot !A_d(\num{d})$.
Nanging iki ateges $d$ nyukupi syarat
definisi~$D$, dadi $d \in D$. Iki
mbantah $d \notin D$. Mula, kanthi
bukti ora langsung, $d \in D$.

Amarga $d \in D$, miturut definisi~$D$,
$\Gamma \Proves \lnot !A_d(\num{d})$.
Ing sisih liyane, amarga $!A_d(x)$
!!{represents}~$D$ ing $\Gamma$,
$\Gamma \Proves !A_d(\num{d})$.
Mula $\Gamma$ ora konsisten.
\end{proof}

\begin{explain}
Teorema sadurunge nuduhake yen ora ana teori
konsisten kang !!{represents} kabeh relasi
!!{decidable} kang uga !!{decidable}.
Kita bakal nuduhake yen $\Th{Q}$ pancen
!!{represents}s kabeh relasi !!{decidable}.
Iki ateges saben teori kang ngemot $\Th{Q}$,
kayata $\Th{PA}$ lan $\Th{TA}$, uga mangkono,
mula uga ora !!{decidable}. (Amarga kabeh
teori kasebut bener ing model baku,
kabeh iku konsisten.)

Asil iki uga bisa digunakake kanggo entuk
versi ringkih saka teorema ora jangkep kapisan.
Saben teori kang !!{axiomatizable} lan
!!{complete} iku !!{decidable}. Mula teori
konsisten kang !!{axiomatizable} lan
!!{represents}s kabeh sifat !!{decidable}
ora bisa !!{complete}.
\end{explain}

\begin{thm}
Yen $\Gamma$ !!{axiomatizable} lan
!!{complete}, teori iku !!{decidable}.
\end{thm}

\begin{proof}
Saben teori kang ora konsisten iku !!{decidable},
amarga teori kaya mangkono ngemot kabeh
!!{sentence}s. Dadi wangsulan kanggo pitakon
``apa $!A \in \Gamma$'' tansah ``ya'',
mula bisa diputusake.

Mula upamane $\Gamma$ konsisten, lan uga
!!{axiomatizable} sarta !!{complete}.
Amarga $\Gamma$ !!{axiomatizable}, teori iku
!!{computably enumerable}. Sebabe, kita bisa
ngenumerasi kabeh !!{derivation}s kang bener
saka aksioma~$\Gamma$ nganggo fungsi komputabel.
Saka !!{derivation} kang bener, kita bisa ngitung
!!{sentence} kang !!{derive}s dening derivasi iku.
Mula ana fungsi komputabel kang ngenumerasi
kabeh teorema~$\Gamma$. Sawijining !!{sentence}
iku teorema saka~$\Gamma$ yen lan mung yen
$\lnot !A$ dudu teorema, amarga $\Gamma$
konsisten lan !!{complete}. Mula kita bisa
mutusake apa $!A \in \Gamma$ kaya mangkene.
Enumerasikna kabeh teorema saka $\Gamma$.
Nalika $!A$ katon ing dhaptar iki, kita ngerti
yen $\Gamma \Proves !A$. Nalika
$\lnot !A$ katon ing dhaptar iki, kita ngerti
yen $\Gamma \Proves/ !A$. Amarga $\Gamma$
!!{complete}, salah siji kedadeyan iki mesthi
pungkasane dumadi, mula prosedur iku pungkasane
menehi wangsulan.
\end{proof}

\begin{cor}
\ollabel{cor:incompleteness}
Yen $\Gamma$ konsisten, !!{axiomatizable},
lan !!{represents} saben sifat !!{decidable},
teori iku ora !!{complete}.
\end{cor}

\begin{proof}
Yen $\Gamma$ !!{complete}, miturut teorema
sadurunge teori iku bakal !!{decidable}
(amarga !!{axiomatizable} lan konsisten).
Nanging amarga $\Gamma$ !!{represents}
saben sifat !!{decidable}, miturut teorema
kapisan teori iku ora !!{decidable}.
\end{proof}

\begin{prob}
Tuduhna yen $\Th{TA} = \Setabs{!A}{\Sat{N}{!A}}$
ora !!{axiomatizable}. Kowé bisa nganggep
yen $\Th{TA}$ nggambarake kabeh sifat
kang bisa diputusake.
\end{prob}

Sawisé kita netepake yen, umpamane, $\Th{Q}$
!!{represents} kabeh sifat !!{decidable},
akibat teorema iku nuduhake yen $\Th{Q}$
mesthi ora jangkep. Nanging buktine ora
menehi conto !!{sentence} mandiri; bukti iku
mung nuduhake yen !!a{sentence} kaya
mangkono mesthi ana. Kanggo nemokake conto,
kita kudu ngaritmetisasi sintaks lan ngetutake
gagasan bukti asli G\"odel. Mesthi wae, kita
uga isih kudu mbuktekake pratelan kapisan:
$\Th{Q}$ pancen !!{represents}s kabeh sifat
!!{decidable}.

Perlu digatekake yen ora saben teori kang
\emph{narik kawigaten} iku ora jangkep utawa
ora bisa diputusake. Ana akeh teori kang
cukup kuwat kanggo njlentrehake fakta
matematika kang narik kawigaten nanging
ora nyukupi syarat asil G\"odel. Tuladhane,
$\Th{Pres} = \Setabs{!A \in \Lang{L_{A^+}}}{\Sat{N}{!A}}$,
yaiku himpunan !!{sentence}s ing basa
aritmetika tanpa~$\times$ kang bener ing
model baku, jangkep lan uga bisa diputusake.
Teori iki diarani aritmetika Presburger,
lan mbuktekake kabeh kabeneran ngenani
bilangan asli kang bisa dirumusake mung
nganggo $\Obj{0}$, $\prime$, lan~$+$.

\end{document}
